% ============================================================================
% Fichier  : partie1/P01_C01_gouvernance_entreprise_si.tex
% Titre    : De la Gouvernance d'Entreprise à la Gouvernance des SI
% Auteur   : MINKA MI NGUIDJOÏ Thierry Emmanuel
% Version  : 1.0 -- octobre 2026
% ============================================================================

\chapter{De la Gouvernance d'Entreprise à la Gouvernance des SI}
\label{chap:gouv-entreprise}
\chaptermark{Gouvernance d'entreprise et des SI}

\epigraph{%
    « Le gouvernement d'entreprise est le système par lequel les sociétés
    sont dirigées et contrôlées. »%
}{Rapport \textsc{Cadbury}, Royaume-Uni, 1992 (notre traduction)}

\niveauChapitre{Fondamental}{%
    Maîtrise du Chapitre~\ref{chap:prologue} (les trois questions réflexes
    et la distinction entre gouvernance et management).%
}

\begin{boxobjectifs}
\begin{itemize}
    \item Définir la gouvernance d'entreprise et ses mécanismes (conseil,
          comités, contrôle).
    \item Expliquer la théorie de l'agence et la théorie des parties
          prenantes.
    \item Retracer l'émergence de la gouvernance des SI (rapport Cadbury,
          loi Sarbanes-Oxley, OHADA).
    \item Définir la gouvernance des SI et la distinguer du management des
          SI.
    \item Identifier les organes de gouvernance des SI (comité stratégique,
          comité de pilotage).
\end{itemize}
\end{boxobjectifs}

\bigskip

% ============================================================================
\section*{Question directrice : pourquoi le conseil ne savait-il rien~?}
% ============================================================================

Le prologue s'est achevé sur une séance extraordinaire du conseil
d'administration de la BANQUE SANAGA. Un administrateur y a posé la
question qui ouvre ce chapitre~: « Nous avons voté 3,2~milliards FCFA pour
ce projet. Comment avons-nous pu ignorer, jusqu'au jour de la panne, qu'il
était en danger~? »

La question n'est pas technique. Elle porte sur la relation entre ceux qui
possèdent l'entreprise, ceux qui la contrôlent et ceux qui la dirigent~:
qui informe qui, qui surveille qui, et avec quels moyens. Cette relation
est l'objet de la \termecle{gouvernance d'entreprise}, discipline bien plus
ancienne que l'informatique. La gouvernance des systèmes d'information n'en
est pas une invention à part~: elle en est le prolongement, devenu
nécessaire le jour où l'activité de l'entreprise a reposé sur ses systèmes.

Ce chapitre suit ce prolongement en cinq temps~: ce qu'est la gouvernance
d'entreprise (section~\ref{sec:c01-gouv-ent}), les théories qui
l'expliquent (section~\ref{sec:c01-theories}), son histoire et son entrée
dans le champ informatique (section~\ref{sec:c01-histoire}), la définition
de la gouvernance des SI (section~\ref{sec:c01-definition}), et enfin ses
organes (section~\ref{sec:c01-organes}), replacés dans le contexte
camerounais (section~\ref{sec:c01-cameroun}).

\separateur

% ============================================================================
\section{La Gouvernance d'Entreprise : Définition et Mécanismes}
\label{sec:c01-gouv-ent}
% ============================================================================

\subsection{Diriger et contrôler}

La définition la plus citée tient en une phrase, celle de l'épigraphe~: la
gouvernance d'entreprise est le système par lequel les sociétés sont
dirigées et contrôlées~\cite{cadbury1992}. Les deux verbes comptent. Une
entreprise dirigée mais non contrôlée dépend entièrement de la qualité de
ses dirigeants. Une entreprise contrôlée mais mal dirigée se conforme sans
avancer. La gouvernance organise l'équilibre entre les deux.

\begin{boxdef}
\termecle{Gouvernance d'entreprise}\index{gouvernance d'entreprise}
(\textit{corporate governance})~: ensemble des relations, des règles et des
mécanismes par lesquels une société est dirigée et contrôlée. Elle répartit
les droits et les responsabilités entre les actionnaires, le conseil
d'administration, la direction générale et les autres parties prenantes,
et fixe les règles de décision et de reddition de comptes.
\end{boxdef}

Les Principes de gouvernance d'entreprise du G20 et de l'\sigle{OCDE},
publiés pour la première fois en 1999 et révisés en dernier lieu en
2023~\cite{oecd2023principles}, constituent la référence internationale.
Ils traitent des droits des actionnaires, de la transparence de
l'information, du rôle des parties prenantes et des responsabilités du
conseil d'administration.

\subsection{Les mécanismes de la gouvernance}

La gouvernance d'entreprise s'appuie sur deux familles de mécanismes. Les
\textbf{mécanismes internes} sont créés par l'entreprise elle-même~; les
\textbf{mécanismes externes} lui sont imposés ou s'exercent de l'extérieur.

\begin{tableaucours}{L{3.6cm}L{4.4cm}L{4.2cm}}{%
    \tableauHeader{Mécanisme} &
    \tableauHeader{Rôle} &
    \tableauHeader{À la BANQUE SANAGA}
}
\multicolumn{3}{l}{\textit{Mécanismes internes}} \\
Assemblée générale & Désigne les administrateurs, approuve les comptes &
Réunie une fois par an \\
Conseil d'administration & Fixe les orientations, contrôle la direction &
A approuvé le projet HORIZON, puis n'en a plus entendu parler \\
Comités du conseil (audit, risques) & Préparent les décisions du conseil
sur des sujets techniques & Comité d'audit sans point sur le projet \\
Contrôle interne et audit interne & Vérifient que les risques sont maîtrisés
& Alerte d'octobre~2025 restée sans suite \\
\multicolumn{3}{l}{\textit{Mécanismes externes}} \\
Commissaire aux comptes & Certifie les comptes & Interroge la fiabilité des
soldes convertis \\
Régulateur (COBAC) & Contrôle la solidité et la conformité & Demande un
rapport circonstancié \\
Marché et clients & Sanctionnent par la réputation et le départ & Files
d'attente, plaintes, presse \\
\end{tableaucours}
\vspace{-0.8em}
\begin{center}{\small\itshape Tableau~1.1. Mécanismes de la gouvernance d'entreprise}\end{center}

Le tableau montre que la BANQUE SANAGA ne manquait d'aucun mécanisme. Ils
existaient tous, mais aucun ne portait sur le système d'information. Le
conseil votait le budget du projet sans en recevoir le suivi~; le comité
d'audit examinait les comptes sans examiner le risque qui allait les
fausser. C'est la situation typique que la gouvernance des SI vient
corriger.

% ============================================================================
\section{Les Théories qui Expliquent la Gouvernance}
\label{sec:c01-theories}
% ============================================================================

\subsection{La théorie de l'agence}

Pourquoi faut-il contrôler les dirigeants~? La réponse la plus influente
vient de Jensen et Meckling, qui décrivent l'entreprise comme un nœud de
contrats, et la relation entre actionnaires et dirigeants comme une
\termecle{relation d'agence}~\cite{jensen1976agency}.

\begin{boxdef}
\termecle{Relation d'agence}\index{relation d'agence} (\textit{agency
relationship})~: contrat par lequel une personne, le \emph{principal},
engage une autre personne, l'\emph{agent}, pour accomplir en son nom une
tâche qui implique de lui déléguer un pouvoir de décision.

\medskip
\termecle{Coûts d'agence}\index{coûts d'agence} (\textit{agency costs})~:
somme des dépenses de surveillance engagées par le principal, des dépenses
d'engagement (de « dédouanement ») engagées par l'agent pour prouver sa
loyauté, et de la perte résiduelle due aux divergences d'intérêts qui
subsistent malgré tout.
\end{boxdef}

La théorie repose sur deux hypothèses. L'agent a ses propres intérêts, qui
ne coïncident pas toujours avec ceux du principal. Et l'agent en sait plus
que le principal sur ce qu'il fait~: c'est l'\termecle{asymétrie
d'information}\index{asymétrie d'information}. La gouvernance est alors
l'ensemble des dispositifs qui réduisent cette asymétrie (rapports,
audits, comités) et alignent les intérêts (rémunération, responsabilité).

\subsection{La chaîne d'agence des systèmes d'information}

L'informatique allonge la chaîne d'agence et accroît l'asymétrie
d'information à chaque maillon. Le conseil délègue au directeur général,
qui délègue au DSI, qui délègue à un intégrateur. À chaque niveau, l'agent
comprend mieux le système que son principal.

\begin{center}
\begin{tikzpicture}[
    maillon/.style = {
        draw=CouleurBordDef, fill=CouleurDef, rounded corners=4pt,
        text width=2.5cm, minimum height=1.0cm, align=center,
        font=\footnotesize\bfseries\color{CouleurPrimaire}
    },
    fleche/.style = {-{Stealth[length=5pt]}, thick, color=CouleurSecondaire}
]
\node[maillon] (ca)  at (0,0)    {Conseil\\ d'administration};
\node[maillon] (dg)  at (3.9,0)  {Directeur\\ général};
\node[maillon] (dsi) at (7.8,0)  {DSI};
\node[maillon] (int) at (11.7,0) {Intégrateur};
\draw[fleche] (ca) -- node[above, font=\scriptsize\itshape]{délègue} (dg);
\draw[fleche] (dg) -- node[above, font=\scriptsize\itshape]{délègue} (dsi);
\draw[fleche] (dsi) -- node[above, font=\scriptsize\itshape]{délègue} (int);
\draw[fleche, color=CouleurAccent] (0,-0.9) -- node[below, align=center,
    font=\scriptsize\itshape\color{CouleurAccent}]{l'information technique
    décroît en remontant la chaîne} (11.7,-0.9);
\end{tikzpicture}
\end{center}
\vspace{-0.4em}
\begin{center}{\small\itshape Figure~1.1. La chaîne d'agence des systèmes
d'information}\end{center}

\begin{boxscenario}{Opération SANAGA. L'agence}
\textbf{Constatation.} Le directeur général fixe seul la date de bascule
en janvier~2026. La DSI l'alerte par une note le 20~février. Le conseil
n'est informé de rien entre juin~2025 et le 9~mars~2026.

\medskip
\textbf{Analyse} (confiance élevée). Le directeur général, agent du
conseil, avait un intérêt propre au respect de la date~: le lancement
de l'offre mobile, déjà annoncé, était l'objectif sur lequel il serait
jugé. Il disposait d'une information (la note de la DSI) que son principal
n'avait pas. Le conseil n'avait mis en place aucun dispositif pour réduire
cette asymétrie~: ni point d'avancement obligatoire, ni accès direct du
comité d'audit à la DSI. La perte résiduelle s'est élevée à
1,9~milliard FCFA.
\end{boxscenario}

\subsection{La théorie des parties prenantes}

La théorie de l'agence ne connaît que deux acteurs~: le principal et
l'agent. Freeman élargit la perspective~: l'entreprise doit tenir compte
de tous les groupes qui affectent ses objectifs ou sont affectés par
eux~\cite{freeman1984stakeholder}. Pour une banque, ce sont aussi les
déposants, les salariés, le régulateur, l'État et les commerçants qui
dépendent de ses paiements.

\begin{boxdef}
\termecle{Partie prenante}\index{partie prenante} (\textit{stakeholder})~:
tout groupe ou individu qui peut affecter la réalisation des objectifs
d'une organisation ou qui est affecté par elle.
\end{boxdef}

Cette théorie a une conséquence directe pour la gouvernance des SI~: la
valeur d'un système d'information ne se mesure pas seulement pour
l'actionnaire. La panne de la BANQUE SANAGA a touché des clients qui ne
pouvaient plus payer leurs salariés, et des commerçants dont les terminaux
de paiement étaient bloqués. COBIT~2019 retient d'ailleurs, comme premier
principe de tout système de gouvernance, la création de valeur pour les
parties prenantes~\cite{isaca2018cobitintro}.

\begin{boiteRemarque}{Et si les dirigeants étaient loyaux~?}
La théorie de l'agence suppose un agent qui sert d'abord ses propres
intérêts. La théorie de l'intendance (\textit{stewardship theory}) part de
l'hypothèse inverse~: un dirigeant peut trouver sa motivation dans la
réussite de l'organisation elle-même~\cite{davis1997stewardship}. Elle
plaide alors pour la confiance et la délégation plutôt que pour la
surveillance. Les deux lectures ne s'excluent pas. Le directeur général de
la BANQUE SANAGA était sans doute sincèrement convaincu de servir la
banque en tenant sa date. Le problème n'était pas sa loyauté, mais
l'absence d'un mécanisme qui confronte sa conviction aux faits techniques.
Une bonne gouvernance ne présume pas la malhonnêteté~: elle organise la
contradiction.
\end{boiteRemarque}

% ============================================================================
\section{Histoire : de Cadbury à la Gouvernance des SI}
\label{sec:c01-histoire}
% ============================================================================

\subsection{Les scandales fondateurs}

La gouvernance d'entreprise s'est construite par réaction à des scandales.
Au Royaume-Uni, la faillite du groupe Maxwell et celle de la banque BCCI,
en 1991, conduisent à la création du comité présidé par Sir Adrian
Cadbury, dont le rapport de 1992 pose les règles de séparation des
pouvoirs et du rôle des comités~\cite{cadbury1992}. Aux États-Unis, les
faillites d'Enron (2001) et de WorldCom (2002) entraînent l'adoption de la
loi Sarbanes-Oxley~\cite{sox2002}.

\subsection{Sarbanes-Oxley~: l'informatique entre dans la gouvernance}

La loi Sarbanes-Oxley est le moment où l'informatique entre dans la
gouvernance d'entreprise, presque sans l'avoir cherché. Sa section~404
impose à la direction des sociétés cotées aux États-Unis d'évaluer chaque
année l'efficacité de son contrôle interne sur l'information financière,
et à l'auditeur externe d'en attester. Or l'information financière est
produite par des systèmes informatiques. Évaluer le contrôle interne,
c'est donc évaluer aussi les contrôles qui garantissent la fiabilité de
ces systèmes~: qui peut modifier un programme, qui peut accéder aux
données, comment les changements sont autorisés.

\begin{boxdef}
\termecle{Contrôles généraux informatiques}\index{contrôles généraux
informatiques} (\textit{IT general controls}, \sigle{ITGC})~: contrôles qui
s'appliquent à l'ensemble de l'environnement informatique et conditionnent
la fiabilité de toutes les applications~: gestion des accès, gestion des
changements, exploitation, développement et acquisition des systèmes.
\end{boxdef}

Le Chapitre~\ref{chap:audit} reviendra sur ces contrôles, que l'auditeur
teste en priorité. Il suffit ici de noter que les défaillances de la
BANQUE SANAGA en sont deux exemples directs~: un changement non autorisé
et des données dont l'exactitude n'était plus garantie.

\subsection{Les étapes de la gouvernance des SI}

\begin{tableaucours}{L{1.6cm}L{10.6cm}}{%
    \tableauHeader{Année} &
    \tableauHeader{Étape}
}
1992 & Rapport Cadbury (Royaume-Uni)~: séparation des pouvoirs, comités
du conseil \\
1994 & Premier rapport King (Afrique du Sud) \\
1996 & Première édition de COBIT, alors centrée sur l'audit et le contrôle
informatiques \\
1997 & Acte uniforme OHADA relatif au droit des sociétés commerciales et du
groupement d'intérêt économique \\
1998 & Création de l'IT Governance Institute (\sigle{ITGI}) \\
2002 & Loi Sarbanes-Oxley (États-Unis) \\
2004 & Weill et Ross définissent la gouvernance des SI par les droits de
décision~\cite{weill2004itgovernance} \\
2008 & Première édition de la norme ISO/IEC~38500 sur la gouvernance des
technologies de l'information \\
2009 & Rapport King~III~: la gouvernance des technologies de l'information
y devient un chapitre à part entière \\
2014 & Acte uniforme OHADA révisé~: comité d'audit obligatoire dans les
sociétés cotées \\
2016 & Rapport King~IV~: principe~12 sur la gouvernance des technologies et
de l'information~\cite{iod2016king4} \\
2018 & COBIT~2019~: système de gouvernance sur mesure, facteurs de
conception \\
2023 & Révision des Principes de gouvernance d'entreprise du G20 et de
l'OCDE \\
\end{tableaucours}
\vspace{-0.8em}
\begin{center}{\small\itshape Tableau~1.2. Repères historiques de la
gouvernance d'entreprise et de la gouvernance des SI}\end{center}

\subsection{Une voix africaine~: les rapports King}

L'Afrique n'est pas restée en marge de ce mouvement. Les rapports King,
publiés par l'Institute of Directors in Southern Africa, figurent parmi
les codes de gouvernance les plus avancés au monde. Le troisième (2009) a
été l'un des premiers codes nationaux à consacrer un chapitre entier à la
gouvernance des technologies de l'information. Le quatrième (2016) en fait
l'un de ses principes~: l'organe de gouvernance doit gouverner la
technologie et l'information de manière à aider l'organisation à définir
et à atteindre ses objectifs stratégiques~\cite{iod2016king4}. COBIT~2019
cite d'ailleurs le rapport King~IV parmi ses références pour l'objectif de
gouvernance EDM01~\cite{isaca2018cobitobjectifs}.

% ============================================================================
\section{Définir la Gouvernance des SI}
\label{sec:c01-definition}
% ============================================================================

\subsection{Quatre définitions de référence}

Il n'existe pas une définition unique, mais quatre définitions qui se
complètent. Chacune met l'accent sur un aspect différent.

\begin{boxdef}
\textbf{COBIT~4.1 (ITGI, 2007, version française de l'AFAI)}~:
« La gouvernance des SI est de la responsabilité des dirigeants et du
conseil d'administration, et elle est constituée des structures et
processus de commandement et de fonctionnement qui conduisent
l'informatique de l'entreprise à soutenir les stratégies et les objectifs
de l'entreprise, et à lui permettre de les élargir. »~\cite{itgi2007cobit41}

\medskip
\textbf{Weill et Ross (2004)}~: spécifier les droits de décision et le
cadre de responsabilité qui encouragent les comportements souhaitables
dans l'usage des technologies de l'information~\cite{weill2004itgovernance}.

\medskip
\textbf{ISO/IEC~38500}~: système par lequel l'utilisation actuelle et
future de l'informatique est dirigée et contrôlée~\cite{iso38500}.

\medskip
\textbf{COBIT~2019}~: activités par lesquelles l'organe de gouvernance
évalue les besoins des parties prenantes, dirige par la priorisation et la
décision, et surveille la performance et la conformité~\cite{isaca2018cobitintro}.
\end{boxdef}

\begin{tableaucours}{L{2.8cm}L{4.4cm}L{5.0cm}}{%
    \tableauHeader{Source} &
    \tableauHeader{Accent principal} &
    \tableauHeader{Question qu'elle pose}
}
COBIT 4.1 & Responsabilité des dirigeants, structures et processus & Qui
porte la gouvernance, et par quels dispositifs~? \\
Weill et Ross & Droits de décision et comportements & Qui décide de quoi,
et quels comportements cela produit-il~? \\
ISO/IEC 38500 & Direction et contrôle de l'usage de l'informatique &
L'usage de l'informatique est-il dirigé et contrôlé~? \\
COBIT 2019 & Évaluer, diriger, surveiller & Les besoins sont-ils évalués,
les priorités fixées, les résultats suivis~? \\
\end{tableaucours}
\vspace{-0.8em}
\begin{center}{\small\itshape Tableau~1.3. Comparaison des définitions de
la gouvernance des SI}\end{center}

\subsection{Une définition de travail pour ce cours}

De ces quatre définitions, ce cours retient une synthèse, qui servira de
référence dans tous les chapitres.

\begin{boxdef}
\termecle{Gouvernance des SI}\index{gouvernance des SI!définition de
travail}~: ensemble des responsabilités, des droits de décision et des
mécanismes par lesquels l'organe de gouvernance d'une organisation
s'assure que l'usage, présent et futur, de l'information et des
technologies \emph{crée de la valeur} pour les parties prenantes,
\emph{maîtrise les risques} qui en découlent et \emph{utilise
judicieusement les ressources} qui lui sont consacrées.
\end{boxdef}

Les trois finalités soulignées (valeur, risque, ressources) seront
développées au Chapitre~\ref{chap:alignement}. Elles donnent déjà un
critère simple~: une décision relève de la gouvernance des SI lorsqu'elle
engage durablement l'une de ces trois finalités pour l'organisation
entière.

\subsection{Structures, processus et mécanismes relationnels}

De Haes et Van Grembergen proposent de décrire la gouvernance des SI d'une
organisation par trois familles de mécanismes~\cite{dehaes2015egit}. La
grille est utile en diagnostic~: une organisation qui n'a que des
structures, sans processus ni relations, a une gouvernance de façade.

\begin{tableaucours}{L{2.6cm}L{4.6cm}L{5.0cm}}{%
    \tableauHeader{Famille} &
    \tableauHeader{Exemples} &
    \tableauHeader{Signal d'absence}
}
Structures & Comité stratégique des SI, comité de pilotage, rôle de DSI
membre de la direction & Personne n'est désigné pour trancher \\
Processus & Planification stratégique des SI, gestion du portefeuille,
suivi de la valeur, audit & Les décisions se prennent au cas par cas,
sans calendrier ni critères \\
Mécanismes relationnels & Formation des administrateurs, communication
entre métiers et DSI, binômes métier-informatique & Les métiers découvrent
les décisions informatiques après coup \\
\end{tableaucours}
\vspace{-0.8em}
\begin{center}{\small\itshape Tableau~1.4. Les trois familles de
mécanismes de gouvernance des SI}\end{center}

% ============================================================================
\section{Gouvernance et Management des SI : Quatre Confusions}
\label{sec:c01-confusions}
% ============================================================================

Le prologue a posé la distinction entre gouvernance et management
(section~\ref{sec:pro-gouv-management}). Dans la pratique, elle est
brouillée par quatre confusions fréquentes. Chacune conduit à une
gouvernance qui n'en a que le nom.

\begin{enumerate}
    \item \textbf{« La gouvernance des SI, c'est l'affaire du DSI. »} Le DSI
          est un acteur central du management des SI. La gouvernance, elle,
          relève de l'organe qui le contrôle~: le conseil et ses comités.
          Confier la gouvernance au DSI, c'est demander à l'agent de se
          surveiller lui-même.
    \item \textbf{« La gouvernance, c'est la conformité. »} Respecter la
          loi et les normes est une exigence de la gouvernance, mais ne
          suffit pas à la définir. Une organisation peut être conforme et
          investir dans des projets sans valeur.
    \item \textbf{« La gouvernance, c'est la sécurité. »} La sécurité de
          l'information est un domaine de risque parmi d'autres. La
          gouvernance arbitre entre la sécurité, la valeur et le coût~;
          elle ne se réduit à aucun des trois.
    \item \textbf{« La gouvernance, c'est de la bureaucratie. »} Une
          gouvernance qui multiplie les comités sans décisions alourdit
          l'organisation. Une bonne gouvernance est proportionnée~: le
          Chapitre~\ref{chap:cobit-conception} montrera comment COBIT~2019
          l'adapte à la taille et aux risques de chaque organisation.
\end{enumerate}

\begin{boxalert}
\textbf{Violation typique.} Après la panne, le conseil d'administration de
la BANQUE SANAGA envisage de « confier la gouvernance des SI » à un DSI
adjoint nouvellement recruté, chargé de produire des rapports mensuels au
directeur général.

\smallskip
\textbf{Règle~:} on ne corrige pas une défaillance de gouvernance en
ajoutant un acteur de management. Il faut un dispositif qui rende compte
\emph{au conseil}, et non à celui dont les décisions doivent être
contrôlées.
\end{boxalert}

% ============================================================================
\section{Les Organes de la Gouvernance des SI}
\label{sec:c01-organes}
% ============================================================================

\subsection{Deux niveaux de comités}

L'ITGI a popularisé une distinction devenue classique entre deux comités.
Le \termecle{comité stratégique des SI}\index{comité stratégique des SI}
(\textit{IT strategy committee}) se situe au niveau du conseil~: il
l'assiste sur les orientations, la valeur et les risques du numérique. Le
\termecle{comité de pilotage des SI}\index{comité de pilotage des SI}
(\textit{IT steering committee}) se situe au niveau de la direction
générale~: il priorise les projets, arbitre les ressources et suit leur
exécution.

\begin{tableaucours}{L{2.8cm}L{2.2cm}L{3.4cm}L{3.8cm}}{%
    \tableauHeader{Organe} &
    \tableauHeader{Niveau} &
    \tableauHeader{Composition type} &
    \tableauHeader{Décisions}
}
Conseil d'administration & Gouvernance & Administrateurs & Approuve la
stratégie numérique, l'appétence au risque, les grands investissements \\
Comité d'audit, comité des risques & Gouvernance & Administrateurs
indépendants & Surveille le contrôle interne et les risques, dont ceux du
SI \\
Comité stratégique des SI & Gouvernance & Administrateurs, DG, DSI,
experts invités & Évalue l'alignement, la valeur et les risques du
portefeuille numérique \\
Comité de pilotage des SI & Management & DG ou DGA, directeurs métiers,
DSI & Priorise les projets, arbitre les ressources, décide des bascules \\
Comité des changements & Management & DSI, exploitation, sécurité,
métiers & Autorise les changements majeurs (Chapitre~\ref{chap:itil-changements}) \\
\end{tableaucours}
\vspace{-0.8em}
\begin{center}{\small\itshape Tableau~1.5. Les organes de la gouvernance et
du management des SI}\end{center}

\subsection{La gouvernance du système de gouvernance}

COBIT~2019 consacre son premier objectif à la mise en place et au maintien
de ce dispositif lui-même. C'est l'objectif qui manquait le plus
visiblement à la BANQUE SANAGA.

\begin{boxobjet}{EDM01 · Assurer la définition et le maintien du cadre de gouvernance}
\textbf{Finalité~:} doter l'organisation d'une approche cohérente de la
gouvernance de l'information et des technologies, intégrée à la
gouvernance d'entreprise, avec des décisions alignées sur sa stratégie.

\textbf{Pratiques~:} EDM01.01 \emph{évaluer} le système de gouvernance
(identifier les parties prenantes et leurs exigences)~; EDM01.02
\emph{diriger} le système de gouvernance (définir les structures, les
principes, les niveaux d'autorité et l'information nécessaire aux
décisions)~; EDM01.03 \emph{surveiller} le système de gouvernance (vérifier
qu'il fonctionne et assure une supervision effective).

\textbf{Sorties~:} principes de gouvernance, modèle de décision et niveaux
d'autorité, rapports sur l'efficacité du système de gouvernance.

\textbf{RACI (extrait)~:} \emph{A} au conseil d'administration, \emph{R} au
directeur général et au comité stratégique des SI, \emph{C} au DSI et à
l'audit interne.

\textbf{Indicateurs~:} part des décisions stratégiques prises par
l'instance prévue~; délai de traitement des alertes remontées au conseil.
\end{boxobjet}
{\small\itshape Intitulés et pratiques traduits par nos soins à partir de
COBIT~2019~\cite{isaca2018cobitobjectifs}.}

\begin{boxprocedure}
\textbf{Un comité de gouvernance qui fonctionne.} Pour qu'un auditeur
puisse établir que le comité stratégique des SI existe autrement que sur
le papier, il doit trouver~:
\begin{itemize}
    \item une charte approuvée par le conseil (mission, composition,
          pouvoirs, fréquence, quorum)~;
    \item un calendrier de réunions respecté, et des procès-verbaux datés
          et signés~;
    \item des décisions inscrites au procès-verbal, avec leur responsable
          et leur échéance~;
    \item un suivi de ces décisions d'une séance à l'autre~;
    \item un compte rendu régulier au conseil d'administration.
\end{itemize}
\end{boxprocedure}

\begin{boxmetier}{Responsable GRC}
\textbf{Rédiger la charte d'un comité stratégique des SI}
\begin{enumerate}
    \item \textbf{Mission~:} assister le conseil dans l'évaluation de la
          stratégie numérique, de la valeur des investissements et des
          risques liés au SI.
    \item \textbf{Composition~:} au moins deux administrateurs, dont l'un
          préside~; le directeur général et le DSI comme membres~; l'audit
          interne comme invité permanent.
    \item \textbf{Pouvoirs~:} accès direct à la DSI et à l'audit interne~;
          droit de demander toute expertise externe.
    \item \textbf{Fréquence~:} au moins trimestrielle, et sur convocation
          en cas de projet majeur en difficulté.
    \item \textbf{Reddition de comptes~:} rapport écrit au conseil après
          chaque séance.
\end{enumerate}
Le point décisif est le deuxième pouvoir~: un comité qui ne reçoit
l'information que par le directeur général reproduit l'asymétrie qu'il
devait réduire.
\end{boxmetier}

% ============================================================================
\section{Le Contexte Camerounais et Africain}
\label{sec:c01-cameroun}
% ============================================================================

\subsection{Des organisations, des cadres différents}

La gouvernance des SI ne s'exerce pas de la même manière dans toutes les
organisations camerounaises, parce que leur gouvernance d'entreprise
relève de cadres juridiques différents.

\begin{tableaucours}{L{3.2cm}L{4.6cm}L{4.4cm}}{%
    \tableauHeader{Organisation} &
    \tableauHeader{Cadre de gouvernance} &
    \tableauHeader{Cas du cours}
}
Société commerciale & Acte uniforme OHADA sur les sociétés commerciales
(révisé en 2014) & KRIBILOG SA, NOUN PAY \\
Établissement de crédit & Droit OHADA et réglementation de la COBAC &
BANQUE SANAGA \\
Entreprise publique & Loi n°~2017/011 du 12~juillet~2017 & · \\
Établissement public & Loi n°~2017/010 du 12~juillet~2017 & ANPP
(opération MEFOU) \\
\end{tableaucours}
\vspace{-0.8em}
\begin{center}{\small\itshape Tableau~1.6. Cadres de gouvernance des
organisations camerounaises}\end{center}

Les lois de 2017 sur les établissements publics et les entreprises
publiques, précisées par le décret n°~2019/320 du 19~juin~2019, organisent
notamment le conseil d'administration et la tutelle de l'État. Le
Chapitre~\ref{chap:droit} en étudiera les dispositions.

% A VERIFIER au Journal officiel de l'OHADA : libellé exact de l'article
% 829-1 de l'AUSCGIE (source utilisée : brochure de présentation de la
% réforme de 2014, PwC/OHADA).
\begin{boxjuris}
\textbf{Acte uniforme OHADA relatif au droit des sociétés commerciales et
du groupement d'intérêt économique, révisé le 30~janvier~2014,
article~829-1}~\cite{ohada2014auscgie}.

\medskip
Selon la présentation de la réforme publiée par l'OHADA, l'Acte uniforme
révisé rend le comité d'audit obligatoire dans les sociétés cotées. Ce
comité, composé exclusivement d'administrateurs non salariés et n'exerçant
pas de mandat de direction générale, assure notamment le suivi de
l'efficacité des systèmes de contrôle interne et de gestion des risques,
et rend compte régulièrement au conseil d'administration.

\medskip
\textbf{Portée.} Le droit OHADA institue ainsi, pour les sociétés cotées,
une instance du conseil chargée de surveiller les risques. Rien ne
l'oblige à examiner les risques du SI, mais rien ne l'en dispense~: le
système d'information porte aujourd'hui l'essentiel du contrôle interne.
Pour les sociétés non cotées, comme la BANQUE SANAGA dans notre cas, la
création d'un tel comité relève des statuts et de la réglementation
sectorielle.
\end{boxjuris}

\subsection{Quatre défis propres au contexte}

\begin{enumerate}
    \item \textbf{La concentration du pouvoir.} Dans de nombreuses
          organisations, le directeur général concentre la décision, et le
          conseil se réunit peu. La chaîne d'agence y est courte en
          apparence, mais l'asymétrie d'information y est forte.
    \item \textbf{La tutelle.} Dans le secteur public, une partie des
          décisions relève du ministère de tutelle. La gouvernance des SI
          doit préciser ce qui se décide au conseil et ce qui remonte à la
          tutelle.
    \item \textbf{La dépendance aux prestataires.} Les grands systèmes sont
          souvent conçus et maintenus par des intégrateurs étrangers. Le
          dernier maillon de la chaîne d'agence échappe alors au droit et
          au contrôle local, d'où l'importance des contrats
          (Chapitre~\ref{chap:itil-niveaux}).
    \item \textbf{La compétence numérique des conseils.} Peu
          d'administrateurs ont une expérience des systèmes d'information.
          La formation des administrateurs est un mécanisme relationnel de
          première importance.
\end{enumerate}

\begin{boxscenario}{Opération MEFOU. Une agence sans comité stratégique}
L'AGENCE NATIONALE DES PAIEMENTS PUBLICS (ANPP), établissement public
installé à Yaoundé, gère le paiement dématérialisé des salaires et des
fournisseurs de l'État. Son schéma directeur informatique, adopté en 2023
pour 14~milliards FCFA, comprend neuf projets. Le conseil d'administration
se réunit deux fois par an et consacre en moyenne vingt minutes au
numérique. Aucun comité ne suit le schéma directeur. Les projets sont
lancés sur décision du directeur général, M.~ESSOMBA, souvent à la demande
du ministère de tutelle.

\medskip
\textbf{Constatation.} À mi-parcours, trois projets ont consommé 71~\% du
budget et aucun n'est achevé. Le conseil l'apprend par une question d'un
administrateur représentant le ministère des Finances.

\medskip
Ce cas reviendra au Chapitre~\ref{chap:alignement} (arbitrage du
portefeuille) et au Chapitre~\ref{chap:audit} (audit de la Chambre des
comptes).
\end{boxscenario}

% ============================================================================
\separateur
% ============================================================================

\begin{boxresume}
\textbf{Ce qu'il faut retenir de ce chapitre~:}
\begin{itemize}
    \item \textbf{Diriger et contrôler}~: la gouvernance d'entreprise est le
          système par lequel les sociétés sont dirigées et contrôlées. Elle
          repose sur des mécanismes internes (conseil, comités, contrôle
          interne) et externes (commissaire aux comptes, régulateur,
          marché).
    \item \textbf{La théorie de l'agence}~: le dirigeant, agent des
          actionnaires, en sait plus qu'eux. La gouvernance réduit cette
          asymétrie d'information. L'informatique allonge la chaîne
          d'agence et l'accroît.
    \item \textbf{Les parties prenantes}~: la valeur d'un SI se mesure pour
          tous ceux qu'il affecte, et non pour le seul actionnaire.
    \item \textbf{L'informatique est entrée dans la gouvernance par le
          contrôle interne}~: la loi Sarbanes-Oxley a fait des contrôles
          généraux informatiques une condition de la fiabilité financière.
    \item \textbf{La gouvernance des SI} est l'ensemble des
          responsabilités, des droits de décision et des mécanismes par
          lesquels l'organe de gouvernance s'assure que l'usage de
          l'information et des technologies crée de la valeur, maîtrise les
          risques et utilise judicieusement les ressources.
    \item \textbf{Deux niveaux de comités}~: le comité stratégique des SI
          relève du conseil, le comité de pilotage relève de la direction
          générale. Ne pas les confondre.
    \item \textbf{Le contexte africain} a ses repères (rapports King, droit
          OHADA, lois camerounaises de 2017) et ses défis (concentration du
          pouvoir, tutelle, dépendance aux prestataires, compétence des
          conseils).
\end{itemize}
\end{boxresume}

\bigskip

\begin{boxexo}[title={\ding{45}\quad Exercice~1.1. Gouvernance ou management~? \niveaubase}]
Classez chacune des décisions suivantes dans la gouvernance ou dans le
management des SI, et indiquez l'organe qui devrait la prendre
(Tableau~1.5).
\begin{enumerate}[(a)]
    \item Fixer la part du budget de la banque consacrée au numérique pour
          les trois prochaines années.
    \item Choisir entre deux éditeurs pour le nouveau logiciel de gestion
          des crédits.
    \item Décider de reporter la bascule d'un système critique.
    \item Fixer le niveau de risque informatique que la banque accepte de
          supporter.
    \item Autoriser l'installation d'un correctif sur les serveurs de
          messagerie.
\end{enumerate}
\end{boxexo}

\begin{boxexo}[title={\ding{45}\quad Exercice~1.2. La chaîne d'agence de KRIBILOG \niveaumoyen}]
La société KRIBILOG SA confie l'exploitation de son logiciel de suivi des
conteneurs à l'infogérant NETPORT SERVICES. Le conseil d'administration ne
reçoit aucun indicateur sur ce service~; le responsable des services
informatiques, M.~MBARGA, reçoit chaque mois un rapport rédigé par
NETPORT SERVICES elle-même.
\begin{enumerate}[(a)]
    \item Dessinez la chaîne d'agence de ce service, sur le modèle de la
          Figure~1.1.
    \item Identifiez, à chaque maillon, l'asymétrie d'information.
    \item Proposez, pour chaque maillon, un mécanisme qui la réduit, en
          précisant s'il s'agit d'une structure, d'un processus ou d'un
          mécanisme relationnel (Tableau~1.4).
\end{enumerate}
\end{boxexo}

\begin{boxexo}[title={\ding{45}\quad Exercice~1.3. Doter l'ANPP d'une gouvernance des SI \niveauexpert}]
Le nouveau président du conseil d'administration de l'ANPP (opération
MEFOU) vous demande une note de trois pages pour doter l'agence d'une
gouvernance des SI.
\begin{enumerate}[(a)]
    \item Rédigez la constatation de départ, à partir des faits de
          l'opération MEFOU, en la séparant de votre analyse.
    \item Proposez les organes à créer, leur niveau (gouvernance ou
          management) et leur composition, en tenant compte de la tutelle
          ministérielle.
    \item Rédigez les cinq articles essentiels de la charte du comité
          stratégique des SI.
    \item Indiquez les trois pièces qu'un auditeur demandera, dans un an,
          pour vérifier que ce dispositif fonctionne réellement.
\end{enumerate}
\end{boxexo}

\bigskip

\begin{boxinfo}
\textbf{Transition.} Ce chapitre a défini la gouvernance des SI et ses
organes à partir de la gouvernance d'entreprise. Le
Chapitre~\ref{chap:defaillances} met ces notions à l'épreuve des faits~:
il analyse de grandes défaillances réelles et documentées, pour montrer
qu'à chaque fois une décision de gouvernance a manqué derrière l'incident
technique.
\end{boxinfo}
